\vspace{1.25in}

\question[15]
Consider the following metapopulation.  Identify \emph{one} subpopulation that is \emph{most likely} to have the highest extinction rate and lowest recolonization rate.  Justify your choice by explaining how patch size, distance and habitat quality interact to affect extinction and recolonization in your chosen patch. There may be more than one correct choice but you only need to explain for one subpopulation. 
\begin{center}
\includegraphics[width=0.65\textwidth]{metapopulation.png}
\end{center}

\newpage

\question[15]
Population size tends to fluctuate around carrying capacity over many generations.  This process may increase the average fitness of the individuals in the population.  Explain how the fluctuation of \textit{N} around \textit{K} may increase the average fitness of the population over many generations. Assume \textit{K} is constant. You may use an illustration to aid (not replace) your explanation. Hint: think about what happens among individuals when \textit{N} is greater than \textit{K}.
